\sectionheading{5.2 প্রোগ্রামিং ভাষার ধারণা}

মানুষ পরস্পরের সঙ্গে যোগাযোগের জন্য বাংলা, ইংরেজি ইত্যাদি ভাষা
ব্যবহার করে। একইভাবে কম্পিউটারকে কোনো নির্দিষ্ট কাজ করানোর জন্যও
নির্দেশ দিতে হয়। কিন্তু মানুষের সাধারণ ভাষায় লেখা নির্দেশ
কম্পিউটার সরাসরি বুঝতে পারে না।

তাই এমন একটি \textbf{নিয়মবদ্ধ ভাষা} প্রয়োজন, যার সাহায্যে
Programmer তার Problem Solving-এর নির্দেশগুলো Program আকারে
লিখতে পারে। এই ভাষাই হলো \textbf{Programming Language}।

%================================================
% Definition
%================================================
\booksubheading{প্রোগ্রামিং ভাষা কাকে বলে?}

\begin{definitionbox}

\textbf{\color{darkTeal}সংজ্ঞা}

কম্পিউটারকে কোনো নির্দিষ্ট কাজ সম্পাদনের নির্দেশ দেওয়ার জন্য
নির্ধারিত \textbf{শব্দ, প্রতীক ও নিয়ম} ব্যবহার করে Program লেখার
ভাষাকে \textbf{প্রোগ্রামিং ভাষা (Programming Language)} বলে।

\end{definitionbox}

\textbf{সহজভাবে মনে রাখবে—}

\begin{center}

\textbf{
Programming Language
\(\;=\;\)
Program লেখার ভাষা
}

\end{center}

অর্থাৎ Programmer তার Problem-এর সমাধানের নির্দেশ
Programming Language-এর সাহায্যে লিখে কম্পিউটারকে কাজ করায়।

%================================================
% Basic relation
%================================================
\booksubheading{Programming Language-এর কাজ}

Programming Language মানুষের Problem Solving Logic এবং
কম্পিউটারের কাজের মধ্যে একটি মাধ্যম হিসেবে কাজ করে।

\begin{center}

\begin{tikzpicture}[node distance=9mm]

\node[databox] (problem)
{\textbf{Problem}\\সমস্যা};

\node[processbox,right=of problem] (logic)
{\textbf{Logic}\\সমাধানের ধাপ};

\node[databox,right=of logic] (program)
{\textbf{Program}\\Language-এ লেখা};

\node[processbox,right=of program] (computer)
{\textbf{Computer}\\কাজ সম্পাদন};

\draw[arrow] (problem) -- (logic);
\draw[arrow] (logic) -- (program);
\draw[arrow] (program) -- (computer);

\end{tikzpicture}

\end{center}

Programmer প্রথমে Problem বুঝে সমাধানের Logic তৈরি করে।
এরপর নির্দিষ্ট Programming Language-এর নিয়ম অনুসরণ করে সেই
Logic Program আকারে লেখা হয়।

%================================================
% Example
%================================================
\begin{examplebox}{উদাহরণ: Screen-এ একটি লেখা প্রদর্শন}

ধরা যাক, কম্পিউটারের Screen-এ \textbf{Hello} প্রদর্শন করতে হবে।

Python Programming Language-এ লেখা যায়—

\begin{center}

{\textenglish{\ttfamily print("Hello")}}

\end{center}

এখানে \textenglish{\ttfamily print()} ব্যবহার করে কম্পিউটারকে
\textbf{Hello প্রদর্শনের নির্দেশ} দেওয়া হয়েছে।

অর্থাৎ—

\begin{center}

\textbf{
নির্দেশ
\(\rightarrow\)
Programming Language-এ Program
\(\rightarrow\)
ফলাফল
}

\end{center}

\end{examplebox}

%================================================
% Elements
%================================================
\booksubheading{Programming Language-এ কী থাকে?}

প্রতিটি Programming Language-এর নিজস্ব \textbf{শব্দ, প্রতীক এবং
লেখার নিয়ম} থাকে। এসব নির্দিষ্ট নিয়ম অনুসরণ করেই Program লেখা হয়।

Programming Language-এর কয়েকটি পরিচিত উপাদান হলো—

\begin{tabularx}{\linewidth}{|>{\bfseries}p{0.20\linewidth}|Y|Y|}
\hline

\rowcolor{softBlue}
\multicolumn{1}{|c|}{\textbf{উপাদান}}
&
\multicolumn{1}{c|}{\textbf{সহজ ধারণা}}
&
\multicolumn{1}{c|}{\textbf{উদাহরণ}}
\\
\hline

Keyword
&
Language-এ আগে থেকে নির্ধারিত বিশেষ শব্দ
&
\textenglish{\ttfamily if, for, while}
\\
\hline

Symbol
&
নির্দিষ্ট কাজ বা অর্থ প্রকাশকারী চিহ্ন
&
\textenglish{\ttfamily +,\ -,\, =}
\\
\hline

Variable
&
Data রাখার জন্য ব্যবহৃত নাম
&
\textenglish{\ttfamily total, age}
\\
\hline

Operator
&
Data-এর ওপর Operation করার চিহ্ন
&
\textenglish{\ttfamily +,\ *,\, /}
\\
\hline

Statement
&
Program-এর কোনো কাজ সম্পাদনের নির্দেশ
&
\textenglish{\ttfamily print("ICT")}
\\
\hline

Syntax
&
Program লেখার সঠিক নিয়ম বা গঠন
&
Language অনুযায়ী নির্ধারিত নিয়ম
\\
\hline

\end{tabularx}

\begin{rememberbox}

\textbf{\color{deepOrange}মনে রাখবে}

প্রতিটি Programming Language-এর নিজস্ব \textbf{Syntax} বা
লেখার নিয়ম রয়েছে। সেই নিয়ম সঠিকভাবে অনুসরণ না করলে Program-এ
Error হতে পারে।

\end{rememberbox}

%================================================
% Example problem
%================================================
\begin{examplebox}{উদাহরণ: দুটি সংখ্যার যোগফল}

Problem: দুটি সংখ্যা \(a\) ও \(b\)-এর যোগফল বের করতে হবে।

মানুষের চিন্তায়—

\[
Sum=a+b
\]

Python-এ—

\begin{center}

\begin{tabular}{@{}l@{}}
\textenglish{\ttfamily a = 10}\\
\textenglish{\ttfamily b = 20}\\
\textenglish{\ttfamily total = a + b}\\
\textenglish{\ttfamily print(total)}
\end{tabular}

\end{center}

Output—

\[
\boxed{30}
\]

এখানে একই Problem-এর সমাধানের Logic একটি Programming Language-এর
নির্দিষ্ট নিয়ম অনুসরণ করে Program আকারে লেখা হয়েছে।

\end{examplebox}

%================================================
% Need
%================================================
\booksubheading{Programming Language কেন প্রয়োজন?}

কম্পিউটারকে কার্যকরভাবে ব্যবহার করার জন্য Programming Language
অত্যন্ত গুরুত্বপূর্ণ। এর সাহায্যে—

\begin{compacttwocolpoints}

\item Problem Solving-এর Logic Program আকারে প্রকাশ করা যায়।

\item কম্পিউটারকে নির্দিষ্ট কাজ করার Instruction দেওয়া যায়।

\item গাণিতিক, যৌক্তিক ও Data Processing-এর কাজ করা যায়।

\item একই ধরনের কাজ স্বয়ংক্রিয়ভাবে বারবার সম্পন্ন করা যায়।

\item Software, Website, Mobile App, Game এবং বিভিন্ন
Automation System তৈরি করা যায়।

\item বড় Problem-কে ছোট অংশে ভাগ করে Program তৈরি করা যায়।

\item Program প্রয়োজন অনুযায়ী পরিবর্তন, উন্নয়ন ও পুনরায় ব্যবহার করা যায়।

\end{compacttwocolpoints}

%================================================
% Human vs programming language
%================================================
\booksubheading{মানুষের ভাষা ও Programming Language-এর পার্থক্য}

মানুষের দৈনন্দিন ভাষা এবং Programming Language—উভয়টিই নির্দেশ বা
তথ্য প্রকাশে ব্যবহৃত হলেও এদের নিয়ম ও ব্যবহার এক নয়।

\begin{tabularx}{\linewidth}{|Y|Y|}
\hline

\rowcolor{softBlue}
\multicolumn{1}{|c|}{\textbf{মানুষের ভাষা}}
&
\multicolumn{1}{c|}{\textbf{Programming Language}}
\\
\hline

মানুষের মধ্যে ভাব ও তথ্য প্রকাশের জন্য ব্যবহৃত হয়
&
কম্পিউটারকে নির্দেশ দিয়ে Program লেখার জন্য ব্যবহৃত হয়
\\
\hline

একই কথা বিভিন্নভাবে প্রকাশ করা যায়
&
নির্দিষ্ট Syntax বা নিয়ম অনুসরণ করে লিখতে হয়
\\
\hline

কিছু ভুল থাকলেও মানুষ Context থেকে অর্থ বুঝতে পারে
&
নিয়মে ভুল হলে Program-এ Error হতে পারে
\\
\hline

শব্দ ও বাক্যের অর্থ অনেক সময় পরিস্থিতিনির্ভর
&
Instruction-এর অর্থ নির্দিষ্ট ও সুনির্দিষ্ট
\\
\hline

উদাহরণ: বাংলা, ইংরেজি
&
উদাহরণ: Python, Java
\\
\hline

\end{tabularx}

%================================================
% Program vs language
%================================================
\booksubheading{Program ও Programming Language-এর পার্থক্য}

\begin{tabularx}{\linewidth}{|Y|Y|}
\hline

\rowcolor{softBlue}
\multicolumn{1}{|c|}{\textbf{Program}}
&
\multicolumn{1}{c|}{\textbf{Programming Language}}
\\
\hline

নির্দিষ্ট কাজ সম্পাদনের জন্য সাজানো Instruction-এর সমষ্টি
&
Program লেখার জন্য ব্যবহৃত নিয়মবদ্ধ ভাষা
\\
\hline

সাধারণত একটি নির্দিষ্ট Problem বা কাজের সমাধান করে
&
বিভিন্ন ধরনের Program তৈরি করতে ব্যবহার করা যায়
\\
\hline

Programming Language ব্যবহার করে Program লেখা হয়
&
Program তৈরির মাধ্যম হিসেবে কাজ করে
\\
\hline

উদাহরণ: দুটি সংখ্যার যোগফল নির্ণয়ের Program
&
উদাহরণ: Python, Java
\\
\hline

\end{tabularx}

\begin{rememberbox}

\textbf{\color{deepOrange}সহজভাবে মনে রাখবে}

\textbf{Programming Language হলো Program লেখার মাধ্যম, আর Program
হলো সেই ভাষায় লেখা কার্যনির্দেশ।}

\end{rememberbox}

%================================================
% Can computer understand?
%================================================
\booksubheading{কম্পিউটার কি সব Programming Language সরাসরি বুঝতে পারে?}

কম্পিউটার মূলত \textbf{Machine Language-এর Instruction} সরাসরি
বুঝতে পারে। Python, Java বা Assembly-এর মতো অন্যান্য ভাষায় লেখা
Program সরাসরি Processor-এর বোধগম্য নয়।

তাই প্রয়োজন অনুযায়ী Program-কে কম্পিউটারের বোধগম্য Instruction-এ
রূপান্তর করতে \textbf{Translator Program} ব্যবহার করা হয়।

\begin{center}

\begin{tikzpicture}[node distance=12mm]

\node[databox] (source)
{\textbf{Source Program}\\Language-এ লেখা};

\node[processbox,right=of source] (translator)
{\textbf{Translator}\\অনুবাদ};

\node[databox,right=of translator] (machine)
{\textbf{Machine Code}\\কম্পিউটারের বোধগম্য};

\draw[arrow] (source) -- (translator);
\draw[arrow] (translator) -- (machine);

\end{tikzpicture}

\end{center}

Compiler, Interpreter এবং Assembler হলো বিভিন্ন ধরনের
Translator Program। এগুলো পরবর্তী অংশে বিস্তারিত আলোচনা করা হবে।

%================================================
% Classification preview
%================================================
\booksubheading{Programming Language-এর স্তরভিত্তিক ধারণা}

কম্পিউটারের Hardware-এর সঙ্গে ঘনিষ্ঠতা এবং মানুষের জন্য
সহজবোধ্যতার ভিত্তিতে Programming Language-কে বিভিন্ন স্তরে
আলোচনা করা হয়।

\begin{tabularx}{\linewidth}{|>{\bfseries}p{0.26\linewidth}|Y|}
\hline

\rowcolor{softBlue}
\multicolumn{1}{|c|}{\textbf{ভাষার ধরন}}
&
\multicolumn{1}{c|}{\textbf{প্রাথমিক ধারণা}}
\\
\hline

Machine Language
&
0 ও 1-এর মাধ্যমে লেখা; কম্পিউটার সরাসরি বুঝতে পারে
\\
\hline

Assembly Language
&
Mnemonic বা প্রতীকী Instruction ব্যবহার করা হয়
\\
\hline

Middle Level Language
&
High Level ও Low Level বৈশিষ্ট্যের সমন্বয় দেখা যায়
\\
\hline

High Level Language
&
মানুষের জন্য তুলনামূলক সহজ ও পাঠযোগ্য
\\
\hline

Fourth Generation Language
&
কম Instruction-এ প্রয়োজনীয় Result প্রকাশে গুরুত্ব দেওয়া হয়
\\
\hline

\end{tabularx}

\begin{rememberbox}

\textbf{\color{deepOrange}খেয়াল করো}

এখানে ভাষাগুলোর শুধু প্রাথমিক ধারণা দেওয়া হয়েছে। Machine Language,
Assembly Language, Middle Level Language, High Level Language এবং
4GL পরবর্তী Topic-গুলোতে বিস্তারিত আলোচনা করা হবে।

\end{rememberbox}

%================================================
% Evolution
%================================================
\booksubheading{Programming Language-এর বিকাশ}

প্রথম দিকের কম্পিউটারে Program লেখা ছিল অনেক কঠিন এবং
Hardware-নির্ভর। সময়ের সঙ্গে Programming Language এমনভাবে
উন্নত হয়েছে যাতে Program লেখা আরও সহজ, বোধগম্য ও কার্যকর হয়।

Programming Language-এর এই বিকাশকে বিভিন্ন
\textbf{Generation} দিয়ে বোঝানো হয়।

\begin{center}

\textbf{
1GL
\(\rightarrow\)
2GL
\(\rightarrow\)
3GL
\(\rightarrow\)
4GL
\(\rightarrow\)
5GL
}

\end{center}

সাধারণভাবে Generation বাড়ার সঙ্গে Language মানুষের Problem
Solving ধারণার কাছাকাছি এসেছে এবং Program লেখা তুলনামূলক সহজ হয়েছে।

%================================================
% Final Remember
%================================================
\begin{rememberbox}

\textbf{\color{deepOrange}মনে রাখবে}

\begin{itemize}

\item Programming Language হলো \textbf{Program লেখার মাধ্যম}।

\item প্রতিটি Programming Language-এর নিজস্ব
\textbf{শব্দ, প্রতীক ও Syntax} থাকে।

\item Program হলো \textbf{নির্দেশনার সমষ্টি}; Programming Language
হলো সেই নির্দেশ লেখার ভাষা।

\item কম্পিউটার মূলত Machine Language সরাসরি বুঝতে পারে।

\item অন্যান্য ভাষায় লেখা Program-কে Machine-এর উপযোগী করতে
Translator Program প্রয়োজন হতে পারে।

\item Programming Language-এর বিভিন্ন স্তর ও Generation রয়েছে।

\end{itemize}

\end{rememberbox}

\small

%================================================
% 5.3 Machine Language
%================================================
